%! Author = markt
%! Date = 3/19/2025

% Preamble
\documentclass[11pt]{article}

% Packages
\usepackage{latexsym,amsfonts,amssymb,amsthm,amsmath, yfonts, cancel, listings, xcolor, graphicx}
\usepackage{color}
\usepackage{listings}

\author{Mark Stanley}

\definecolor{codegreen}{rgb}{0,0.6,0}
\definecolor{codegray}{rgb}{0.5,0.5,0.5}
\definecolor{codepurple}{rgb}{0.58,0,0.82}
\definecolor{backcolour}{rgb}{0.95,0.95,0.92}
\lstdefinestyle{mystyle}{
    backgroundcolor=\color{backcolour},
    commentstyle=\color{codegreen},
    keywordstyle=\color{magenta},
    numberstyle=\tiny\color{codegray},
    stringstyle=\color{codepurple},
    basicstyle=\ttfamily\footnotesize,
    breakatwhitespace=false,
    breaklines=true,
    captionpos=b,
    keepspaces=true,
    numbers=left,
    numbersep=5pt,
    showspaces=false,
    showstringspaces=false,
    showtabs=false,
    tabsize=2
}

\title{Math 443 Session 22}

% Document
\begin{document}
    \lstset{style=mystyle}

    \maketitle

    \textbf{Least Square Approximation}
    We can fit a polynomial of degree $n$ to $y(t)=\alpha_0 + \alpha_1 t + \cdots + \alpha_{n} t^n$

    We can write this as a matrix equation $Ax=y$, where $A$ is the Vandermonde matrix.

    \[A=\begin{bmatrix}
            1 & t_1 & t_1^2 & \cdots & t_1^n \\
            1 & t_2 & t_2^2 & \cdots & t_2^n \\
            \vdots & \vdots & \vdots & \ddots & \vdots \\
            1 & t_m & t_m^2 & \cdots & t_m^n
        \end{bmatrix}\]

    \textbf{Theorem}
    Let $t_1, \ldots, t_{n+1}$ be distinct sample points. Then, for any prescribed data $y_1, \ldots, y_{n+1}$, there
    exists a unique polynomial $p(t)$ of degree at most $n$ such that $p(t_i)=y_i, \quad \forall i=1, \ldots, n+1$


    \begin{lstlisting}[language=Python]
    def python_code():
        pass
    \end{lstlisting}

\end{document}